% system-model.tex -- EVM Executor Paper (BNR-2026-008)

\section{System Model}

\subsection{Target Architecture}

The Basis Network enterprise zkEVM L2 operates a per-enterprise chain
architecture where each client enterprise runs an independent L2 instance.
Each instance consists of:

\begin{itemize}
  \item A \textbf{Go-based sequencer/executor node} that receives transactions,
    executes them against a Poseidon-based sparse Merkle tree (SMT), and
    produces structured execution traces.
  \item A \textbf{Rust-based ZK prover} that consumes execution traces via
    gRPC, generates PLONK validity proofs, and returns proof artifacts.
  \item A \textbf{proof submitter} that posts validity proofs and state roots
    to the Basis Network L1 (Avalanche Subnet-EVM, Chain~ID 43199).
\end{itemize}

This architecture separates execution (Go) from proving (Rust), matching the
established pattern observed in all five production zkEVM systems surveyed
in Section~II.

\subsection{EVM Execution Requirements}

The execution engine must satisfy the following formal requirements:

\begin{enumerate}
  \item \textbf{Opcode completeness}: For every Cancun EVM opcode $o \in O_{cancun}$
    and valid input state $\sigma$, the execution result must be identical to
    the reference Geth implementation:
    $\forall o \in O_{cancun}, \sigma: \text{exec}_{basis}(o, \sigma) = \text{exec}_{geth}(o, \sigma)$.

  \item \textbf{Trace completeness}: For every executed transaction $tx$, the
    trace $T(tx)$ must contain all state-modifying operations:
    $T(tx) \supseteq \{(op, addr, slot, val) \mid op \in \{$SLOAD, SSTORE,
    BALANCE$\_$CHANGE, NONCE$\_$CHANGE, LOG, CREATE$\}\}$.

  \item \textbf{Throughput bound}: With selective ZK tracing enabled,
    throughput must exceed 1{,}000 transactions per second for simple
    transfers (21{,}000 gas each).
\end{enumerate}

\subsection{Geth Module Dependency Analysis}

We analyze go-ethereum's module structure to identify the minimal set of
packages required for standalone EVM execution with tracing.

\begin{table}[htbp]
\caption{Go-Ethereum Modules Required for Standalone Execution}
\label{tab:geth-modules}
\centering
\begin{tabular}{llr}
\toprule
\textbf{Module} & \textbf{Purpose} & \textbf{LOC} \\
\midrule
\texttt{core/vm/} & EVM interpreter, opcodes, stack, memory & ${\sim}$8{,}000 \\
\texttt{core/vm/runtime/} & Standalone execution helpers & ${\sim}$500 \\
\texttt{core/state/} & StateDB, stateObject, journal & ${\sim}$5{,}000 \\
\texttt{core/types/} & Transaction, Block, Receipt, Log & ${\sim}$4{,}000 \\
\texttt{ethdb/} & Database interface (LevelDB, memorydb) & ${\sim}$1{,}000 \\
\texttt{params/} & Chain config, protocol params & ${\sim}$2{,}000 \\
\texttt{crypto/} & Keccak hashing, secp256k1 signing & ${\sim}$1{,}500 \\
\texttt{common/} & Address, Hash, Big types & ${\sim}$2{,}000 \\
\texttt{core/tracing/} & Tracing hooks interface & ${\sim}$500 \\
\midrule
\textbf{Subtotal (required)} & & ${\sim}$\textbf{24{,}500} \\
\bottomrule
\end{tabular}
\end{table}

\begin{table}[htbp]
\caption{Go-Ethereum Modules Not Required}
\label{tab:geth-excluded}
\centering
\begin{tabular}{llrl}
\toprule
\textbf{Module} & \textbf{Purpose} & \textbf{LOC} & \textbf{Exclusion Reason} \\
\midrule
\texttt{p2p/} & Peer networking & ${\sim}$15{,}000 & Custom L2 networking \\
\texttt{eth/} & Protocol handler & ${\sim}$20{,}000 & Full node protocol \\
\texttt{les/} & Light client & ${\sim}$10{,}000 & Not applicable \\
\texttt{consensus/} & PoW/PoS consensus & ${\sim}$8{,}000 & L2 uses sequencer \\
\texttt{miner/} & Block mining & ${\sim}$3{,}000 & Not applicable \\
\texttt{cmd/} & CLI tools & ${\sim}$10{,}000 & Custom node binary \\
\texttt{internal/} & Internal APIs & ${\sim}$5{,}000 & Custom APIs \\
\midrule
\textbf{Subtotal (excluded)} & & ${\sim}$\textbf{71{,}000} & \\
\bottomrule
\end{tabular}
\end{table}

The estimated fork ratio is approximately 24{,}500 / 95{,}500 $\approx$ 25\%
of core Go code. In practice, the preferred approach imports go-ethereum as a
Go module dependency rather than forking the repository, using only the
required packages (Table~\ref{tab:geth-modules}).

\subsection{Fork Strategy Decision Framework}

Two strategies exist for integrating Geth's EVM:

\textbf{Strategy~A (Import as Go module):}
\begin{lstlisting}[language=Go]
import (
  "github.com/ethereum/go-ethereum/core/vm"
  "github.com/ethereum/go-ethereum/core/state"
)
\end{lstlisting}
This provides automatic upstream updates, zero maintenance burden, and clean
dependency management, but limits modification to public APIs.

\textbf{Strategy~B (Repository fork):}
Full control over all internals, including the ability to modify StateDB,
add custom opcodes, and alter execution semantics. The cost is merge
conflicts on upstream updates and ongoing maintenance overhead.

Based on the analysis in this paper, we recommend Strategy~A. The
\texttt{vm.StateDB} interface is sufficiently abstract to accommodate a
custom Poseidon SMT implementation, and the \texttt{core/tracing} hooks
API provides the trace capture capabilities needed for ZK witness generation.
